\documentclass[10pt]{article}
\usepackage[margin=20mm]{geometry}
\usepackage{amsmath,amssymb,fontspec,unicode-math,xurl,hyperref}
\IfFontExistsTF{DejaVu Serif}{\setmainfont{DejaVu Serif}}{\setmainfont{Latin Modern Roman}}
\setmathfont{latinmodern-math.otf}
\setlength{\parindent}{0pt}\setlength{\parskip}{6pt}\sloppy
\begin{document}
Abstract: The paper studies whether gradient flow learns nearly optimally robust classifiers for an orthonormal Gaussian mixture. Two demonstrated errors affect the central argument. Lemma 17 converts a cosine deficit into a directional error without a square root; an exact balanced two-neuron example violates its bound by a factor of approximately \(30.9\). The proof of the main uniform approximation theorem explicitly invokes this lemma. Separately, the robustness bridge treats a degree-\(p\) function as degree one. Both are classified Major because they invalidate derivations on the core dependency path. Neither counterexample proves that an actual sampled gradient-flow trajectory fails to become robust. A corrected directional bound and a replacement homogeneity argument delimit what remains plausible.\par
Author: AI-assisted technical review\par
Date:\par
Title: Initial Scientific Review of
Gradient Flow Provably Learns Robust Classifiers for Orthonormal GMMs\par
\# Identity, scope, and result\par
This is the initial review of paper 0185, I25-43554, by Hancheng Min and René Vidal, in the ICML 2025 proceedings, PMLR 267, pp. 44292–44350. The reviewed artifact is the staged 59-page proceedings PDF, identified by SHA-256\par
> c5fce6403962d90a387f98ae9e7043d1b9a566d25c14e29366f8d04750b6fd80.\par
All locations below use its one-based PDF page numbers. The source manifest and original PDF identity were checked; no replacement version was downloaded. No paper-specific code was identified in the staged preparation. Access to a separate OpenReview attachment remains unresolved, rather than evidence that no attachment exists.\par
Assessment: Major. The findings concern explicit false inequalities and an invalid core proof step. Scientific adjudication remains pending; no human review or author adoption is recorded. The auxiliary agent examined the main text and selected appendices. The parent checked the decisive original equations, source images on pp. 4, 22, 35 and 59, the main-theorem dependency on pp. 6 and 36, and independently recomputed the exact counterexample. This is not a certification of every dynamical lemma in the 59-page paper.\par
\# Contribution and dependency structure\par
For orthonormal cluster centers \(\mu_1,\ldots,\mu_K\), the paper considers a balanced Gaussian-mixture training set and the network \[f^{(p)}(x;\theta)=\sum_j v_j
 \frac{[\langle x,w_j\rangle]_+^p}{\|w_j\|^{p-1}},\qquad p>2,\] with small, balanced, nondegenerate initialization. Under balance \(|v_j|=\|w_j\|\), each neuron contributes \(\operatorname{sign}(v_j)\|w_j\|^2[\langle x,u_j\rangle]_+^p\), where \(u_j=w_j/\|w_j\|\). The target classifier is \[F^{(p)}(x)=\sum_{k\le K_1}[\langle x,\mu_k\rangle]_+^p
 -\sum_{k>K_1}[\langle x,\mu_k\rangle]_+^p.\] Theorem 1, Eq. (12), asserts uniform sphere approximation of order \(\widetilde O(\alpha_0^2)\) over a training interval, under its dimension, sample-size, initialization and variance assumptions. Proposition 3 converts suitable sphere approximation into near-optimal adversarial robustness. Lemma 17 is the stated conversion from directional alignment and cluster weight control to the uniform approximation. These dependencies make both problems central despite their detailed derivations appearing in appendices.\par
\# Problem 1: a false angular-to-uniform error bound\par
Rating: Major. Locations: Lemma 17, p. 35 and pp. 58–59; Theorem 1, p. 6, Eq. (12); restated theorem, p. 35, Eq. (G.5); explicit application of Lemma 17, p. 36.\par
\#\# Claim and assumptions\par
The lemma assumes cosine alignment within a deficit \(\delta\), per-cluster squared-weight sums within \(\nu\) of one, and total void-neuron squared weight at most \(\zeta\). It concludes \[\sup_{\|x\|=1}|f^{(p)}(x;\theta)-F^{(p)}(x)|
 \le K(1+\nu)(2^p-1)2\delta+K\nu+\zeta. \tag{R1}\] The printed alignment hypothesis uses maxima; the following counterexample satisfies the stronger, intended condition for every neuron, so correcting that quantifier does not remove the violation.\par
\#\# Exact admissible counterexample\par
Set \(D=3\), \(K=2\), \(K_1=K_2=1\), \(p=3\), and use two neurons. Let \[t=\frac1{1000},\quad a=\frac{1-t^2}{1+t^2},\quad
 b=\frac{2t}{1+t^2},\quad
 \mu_1=e_1,\quad\mu_2=e_2,\] \[w_1=(a,0,b),\quad v_1=1,\qquad w_2=e_2,\quad v_2=-1.\] Take \(\mathcal N_1=\{1\}\), \(\mathcal N_2=\{2\}\), no void neurons, \(\delta=1-a=2/1000001\), and \(\nu=\zeta=t^4>0\). Both neurons have unit norm and satisfy balance; each cluster has squared-weight sum exactly one. Their corresponding cosines are \(a\) and one, respectively. The centers are orthonormal. At the unit vector \(x=(3/5,0,4/5)\), \[|f^{(3)}(x)-F^{(3)}(x)|
 =\left(\frac{3a+4b}{5}\right)^3-\left(\frac35\right)^3
 =\frac{216414079424215946}{125000375000375000125}
 \approx0.001731307441466209.\] In contrast, the right side of (R1) is \[28(1+t^4)(1-a)+3t^4
 =\frac{56000003000059}{1000001000000000000}
 \approx0.000055999947000112.\] The ratio is approximately \(30.9162\). Both an auxiliary exact-rational diagnostic and the parent’s separately written Fraction calculation reproduced this violation. No author code was executed.\par
\#\# Failure mechanism and bounded consequence\par
The displayed proof on p. 59 replaces \(\|u_j-\mu_k\|\) by \(2(1-c_{kj})\). The correct identity for unit vectors is \[\|u_j-\mu_k\|^2=2(1-c_{kj}),\qquad
 \|u_j-\mu_k\|=\sqrt{2(1-c_{kj})}.\] The example’s output difference is \(216t/125+O(t^2)\), whereas (R1) is \(56t^2+O(t^4)\). This is a rate error, not rounding or a loose constant.\par
The theorem proof on p. 36 explicitly uses Lemma 17 after obtaining cosine deficits of order \(\widetilde O(\alpha^2)\). Thus its stated deduction of the core uniform \(\widetilde O(\alpha_0^2)\) rate is invalid. This is sufficient for Major under a core-derivation standard. The example is a configuration satisfying the lemma, not a constructed sampled gradient-flow trajectory satisfying every assumption of Theorem 1. It therefore does not independently disprove that theorem’s conclusion or the paper’s qualitative learning mechanism. Extra dynamical cancellation or a stronger alignment estimate could in principle recover a sharper rate, but it is not supplied by this lemma.\par
\#\# Verified local correction\par
Since \(z\mapsto[z]_+^p\) is \(p\)-Lipschitz on \([-1,1]\), a direct replacement, under the all-neuron alignment hypothesis, is \[\sup_{\|x\|=1}|f^{(p)}(x)-F^{(p)}(x)|
 \le K(1+\nu)p\sqrt{2\delta}+K\nu+\zeta. \tag{R2}\] This follows by summing each neuron’s directional perturbation and then bounding the cluster mass errors. The void contribution is at most \(\zeta\). Retaining the cluster-count dependence, angular and mass errors of order \(\widetilde O(\alpha_0^2)\) make this generic route certify a sphere-error upper bound \(\widetilde O(K\alpha_0)\) for fixed \(p\). At fixed \(K,p\), this is \(\widetilde O(\alpha_0)\), whereas the printed route claims \(\widetilde O(\alpha_0^2)\). Conditional on all other dynamical hypotheses and a sufficiently small tolerance for the repaired robustness bridge, the corresponding certified radius deficit is \(\widetilde O(K^{1/p}\alpha_0^{1/p})\), or \(\widetilde O(\alpha_0^{1/p})\) at fixed \(K,p\). The paper’s tilde convention does not hide polynomial dependence on \(K\). These are upper bounds supplied by the replacement route, not lower bounds excluding a faster rate on actual trajectories and not a replacement theorem for the full dynamics.\par
\# Problem 2: incorrect homogeneity in the robustness bridge\par
Rating: Major. Locations: Proposition 3 and the paragraph following it, p. 5; Appendix D, p. 22, especially Eq. (D.2).\par
\#\# Demonstrated invalid step\par
Proposition 3 assumes \(f(\gamma x)=\gamma f(x)\). Its proof states that the same holds for \(F^{(p)}\), and uses a sphere discrepancy at most \(\nu\) to infer an off-sphere discrepancy at most \(\|x\|\nu\). The main text also calls the actual network degree one in its input. From the definitions, however, \[f^{(p)}(\gamma x;\theta)=\gamma^p f^{(p)}(x;\theta),\qquad
 F^{(p)}(\gamma x)=\gamma^p F^{(p)}(x).\] The denominator in the network depends on the weights, not on the input, and does not change this input homogeneity.\par
There is a direct counterexample to the proof step even under Proposition 3’s stated degree-one premise. Define \(f(x)=\|x\|^{1-p}F^{(p)}(x)\) for \(x\ne0\) and \(f(0)=0\). It is degree one and has zero sphere discrepancy from \(F^{(p)}\). For \(p=3\), \(\mu_1=e_1,\mu_2=e_2\), and \(x=2e_1\), its value is 2 while \(F^{(3)}(x)=8\). The discrepancy is 6. Choosing \(\nu=0.001<(\sqrt2/8)^3\) satisfies the proposition’s tolerance assumption, while the claimed off-sphere bound is only \(0.002\).\par
This disproves Eq. (D.2)’s inference and shows that Proposition 3 as stated cannot directly be applied to the degree-\(p\) network. It does not by itself disprove Proposition 3’s robustness conclusion, because classification signs are invariant under positive radial rescaling.\par
\#\# A replacement bridge and its limits\par
A local repair is possible. For a positively homogeneous classifier of any positive degree, let \(g(z)=\|z\|^p f(z/\|z\|)\), with \(g(0)=0\). Then \(g\) has the same sign as \(f\) and the same unit-sphere values. If \(\sup_{\|u\|=1}|cf(u)-F^{(p)}(u)|\le\nu\) for \(c>0\), the correct extension is \[|cg(z)-F^{(p)}(z)|\le\|z\|^p\nu.\] The distance definition can be handled by a limiting argument, or by attainment: the relevant supremum norm is coercive in \(c\) for nonzero bounded \(f\), and \(\|F^{(p)}\|_\infty=1>\nu\) excludes a zero minimizing scale.\par
For completeness, put \(q=\nu^{1/p}\) and \(r=(\sqrt2-8q)/2\), so \(0<q\le\sqrt2/8\). Write a test sample as \(\mu_k+\varepsilon\) and restrict to \[\max_l|\langle\varepsilon,\mu_l\rangle|\le q/K,\qquad
 \|\varepsilon\|\le\sqrt{17/2}-1.\] For any \(\|d\|\le1\), \(z=\mu_k+\varepsilon+rd\) satisfies \(\|z\|\le R=\sqrt{17/2}+\sqrt2/2\). The correct-class coordinate minus the opposite-class \(\ell_p\) norm is at least \[1-q-r\sqrt2=(4\sqrt2-1)q>Rq.\] Writing \(J\) for the opposite-label clusters, \(m=|J|\le K-1\), and \(\varepsilon_l=\langle\varepsilon,\mu_l\rangle\), the two bounds used here are \[|\varepsilon_k|+\|\varepsilon_J\|_p
 \le (1+m^{1/p})q/K\le q,\qquad
 |d_k|+\|d_J\|_2\le\sqrt2.\] The second uses orthonormality and \(\|d\|_2\le1\); the first uses \(p>1\) and both classes being represented. For nonnegative \(A\ge B\) and \(p>1\), \(A^p-B^p\ge(A-B)^p\). Thus the correct-class margin exceeds \(R^p\nu\), which dominates the approximation error. The same argument applies to either class. Combining this deterministic argument with the Gaussian projection and norm tail bounds gives the stated probability form \[1-2K\exp\!\left(-\frac{D\nu^{2/p}}{2K^2\alpha^2}\right)
 -4\exp\!\left(-\frac{3}{8\alpha^2}\right).\] The constants satisfy \(4\sqrt2-1\approx4.65685>R\approx3.62258\) and \((\sqrt{17/2}-1)^2>3\), independently checked numerically. The replacement also avoids the original proof’s unjustified norm-three restriction. The norm-tail bound needs no small-\(\alpha\) condition: for \(\varepsilon\sim N(0,\alpha^2I/D)\), \[{\mathbb E}\exp\!\left(\frac{\|\varepsilon\|^2}{4\alpha^2}\right)
 =(1-1/(2D))^{-D/2}\le\sqrt2.\] Markov’s inequality and \((\sqrt{17/2}-1)^2>3\) give the displayed weaker \(4\exp[-3/(8\alpha^2)]\) tail for every \(\alpha>0\) and integer \(D\ge1\). The deterministic margin is uniformly at least \(\{(4\sqrt2-1)^p-R^p\}\nu>0\) before division by the positive scale. Attacked inputs stay in an annulus bounded away from zero, so transferring this margin back to any positive homogeneity degree is valid. For arbitrary discontinuous classifiers, use an infimum over attacks; the actual network is continuous.\par
This repair explains why the robustness conclusion can survive the homogeneity error. It does not retroactively validate the printed derivation and does not repair Problem 1’s main approximation rate. Major here records a demonstrated invalid core bridge, with unusually strong counterevidence against claiming the robustness principle itself false.\par
\# Conclusion and verification status\par
Scientific adjudication remains pending; human reviews are empty. The issue grades below are unchanged. Adoption has not been requested or verified.\par
Issue assessment | Repair verification | Adoption | Remaining work\par
Issue assessment | Repair verification | Adoption | Remaining work\par
P1: Major; Lemma 17 loses a square root. | Local \(\sqrt{\delta}\) bound checked with explicit \(K\) dependence. | Not requested. | Re-establish the main dynamical approximation rate under its full assumptions.\par
P2: Major; invalid homogeneity step on the robustness proof path. | Degree-p radial extension, deterministic margin and Gaussian-tail bridge checked locally. | Not requested. | Integrate the conditional bridge with a valid trajectory theorem; P1 remains unresolved.\par
9 Hancheng Min and Rene Vidal. Gradient Flow Provably Learns Robust Classifiers for Orthonormal GMMs. Proceedings of Machine Learning Research 267 (ICML 2025), 44292–44350. https://proceedings.mlr.press/v267/min25c.html. Reviewed proceedings version.\par
\# Coverage, evidence, and unresolved work\par
\#\#\#\# Recorded calculation and evidence identity.\par
The parent-executed reviewer script `parent\_check.py` and its completed `parent\_check\_result.json` are preserved under . It records the exact Lemma 17 witness, degree-mismatch diagnostic, and local robustness-margin constants. These are reviewer checks, with no author-code execution. Native auxiliary roles recorded for this paper: (completed), (completed; parent integrated). Exposed runtime identities and scope-specific source-image observations are retained in the accompanying provenance record; this endmatter pass does not claim a new scientific execution.\par
> Script SHA-256: .
> Result SHA-256: .\par
The auxiliary review covered the main text, Appendix A, the midpoint and Bayes-classifier arguments in Appendix C, Appendix D, balance and homogeneity in E.1, statements in F.1 and G.1, the main proof in G.2, and Lemma 17’s full proof. It found no demonstrated issue in the inspected midpoint-counting bound or the Bayes log-sum-exp calculation. Most remaining detailed dynamical proofs in F and H, the full good-event concentration derivation, and empirical illustrations were not exhaustively rederived.\par
Durable evidence includes `auxiliary/scientific\_notes.md`, its coverage and hash manifest, the exact `check\_lemma17.py` and result, and the separate parent `parent\_check.py` and result. Source-page images are preserved and were actually opened. Machine-readable coverage distinguishes auxiliary inspection from the parent’s decisive checks. A separate challenge by independently reproduced the exact witness, verified the radial repair and its Gaussian tail, and identified the necessary explicit \(K\) dependence above. The parent read its full analysis and integrated these qualifications; the additional challenge execution is attributed to that auxiliary, not to a second parent execution. No experiment, training run, author-code execution, external dataset or weight download was performed. No full later-version or erratum search was completed.\par
\# Production and adjudication status\par
The final English TeX retains the frozen mathematical-journal preamble. The initial Windows-only pass lacked an available TeX compiler. Recovery build inputs, output hashes, and actual final-report visual checks are recorded in `provenance.json`; original-source page inspections remain separately attributed. Scientific adjudication remains pending, with no human review or author adoption claimed.\par
\end{document}
